\answer{ex:02:1}{길이 $3.2$~mm, 간격 $25\um$인 열: 검출기 $\approx129$개. 폭 $v\,\tau\ln2=1.7$~mm인 띠, 즉 검출기 $\approx69$개가 발화하는데, 열의 절반이 넘는다. 방향은 띠의 중심이다.}
\answer{ex:02:2}{창은 $-\tau\ln\frac{w_2}{\theta-w_1}\le\delta\le\tau\ln\frac{w_1}{\theta-w_2}$, 즉 $[-0.235,\,0.178]\ms$; 중심 $c=\frac\tau2\ln\frac{w_1(\theta-w_1)}{w_2(\theta-w_2)}=-28$~\textmu s. 방향이 강한 귀 쪽으로 $28$~\textmu s 밀린다. 분해능 간격의 거의 세 배이다.}
\answer{ex:02:3}{$s_iw_{ij}s_j\ge0$인 $s_i$가 필요하며, 이런 $s_i$는 순환이 짝수일 때 그리고 그때에만 존재한다. 상호 억제는 귀착된다(지속 진동 없음). \nt{GLUT}--\nt{GABA} 루프는 귀착되지 않는다(진동할 수 있음).}
\answer{ex:02:4}{뉴런 여섯 개 $g_k=[x+y+z\ge k]$와 $\bar g_k=[\bar x+\bar y+\bar z\ge4-k]$, $k=1,2,3$; $C=g_2$, $\bar C=\bar g_2$, $S=[g_1+\bar g_2+g_3\ge2]$, $\bar S=[\bar g_1+g_2+\bar g_3\ge2]$: 뉴런 여덟 개. AND와 OR로는 $12$개.}
\answer{ex:02:5}{$T_{\min}(0.6)=0.718\,\tau$; $I_c=0.225$, 문턱 근처에서 $T_{\min}\propto(I-I_c)^{-1/2}$.}
\answer{ex:02:6}{(a)~고리 $500$개, 뉴런 $5\cdot10^4$개; 산포 $4.5\ms$($0.45\,\%$), 상대 오차 $\propto T^{-1/2}$. (b)~모든 고리에 공통인, 산포 $5\,\%$의 무작위 지연 배율.}
\answer{ex:02:7}{$\tau_F\ln2=0.69\,\text{s}$마다 갱신: 뉴런당 초당 스파이크 $4.3$개 대 $20$개로 $4.6$배 경제적이다(스파이크당 $10^{-10}$~J일 때 $4\cdot10^{-7}$ 대 $2\cdot10^{-6}$~J). 플립플롭은 비트를 잃는다. 커패시터는 침묵이 갱신 뒤 $0.49\,\text{s}$ 이내에 시작되면 비트를 지킨다(확률 $0.71$).}
